\section{Calibration procedures and results}
\label{app:calibration-details}

Every calibrator is fitted on the dedicated calibration holdout inside the corresponding outer
training split. Table~\ref{tab:calibration-map} summarizes how each form of uncertainty is
calibrated. Figures~\ref{fig:binary-dataset-scores}, \ref{fig:appendix-native-binary-nll}, and
\ref{fig:appendix-native-binary-ece} show Brier score, binary NLL, and ECE by dataset. Figures
\ref{fig:regression-dataset-scores}, \ref{fig:appendix-native-regression-nll}, and
\ref{fig:appendix-native-regression-coverage-error} show interval score, Gaussian NLL, and coverage
error for regression.

\begin{table*}[!htbp]
\centering
\caption{Calibration applied to each pipeline. A supplied scale is a standard deviation on the
target used to train the model.}
\label{tab:calibration-map}
\footnotesize
\begin{tabularx}{\textwidth}{@{}>{\raggedright\arraybackslash}p{0.29\textwidth}
  >{\raggedright\arraybackslash}p{0.31\textwidth}>{\raggedright\arraybackslash}X@{}}
\toprule
Pipelines & Supplied uncertainty & Calibration applied \\
\midrule
All eleven binary pipelines & Marginal event probability & One scalar temperature \\
LightGBM; deterministic MLP & Point prediction only & Constant Gaussian residual scale \\
MC dropout; SWAG; deep ensemble; ensemble MC dropout; MultiSWAG
  & Stochastic-pass, posterior-draw, or ensemble-member spread
  & Additive observation-noise variance \\
Laplace; post-hoc VBLL; MoLA; ensemble post-hoc VBLL
  & Total predictive standard deviation & Standard-deviation multiplier \\
\bottomrule
\end{tabularx}
\end{table*}

For binary predictions, we choose the temperature that minimizes NLL on the calibration set and
divide the probability logits by this value. If the calibration labels cannot support a finite
temperature estimate, we set the temperature to one and keep the original probabilities.

For regression pipelines with only a point prediction, the fitted constant scale is the
root-mean-square calibration residual. For pipelines with a predictive standard deviation, the
scale multiplier is the root-mean-square residual after division by that standard deviation. For
pipelines supplying epistemic spread, we choose a non-negative observation-noise variance by
minimizing Gaussian NLL on the calibration set and add it to the supplied variance. Central 90\%
intervals extend 1.645 adjusted standard deviations on either side of the mean.

The four log-trained endpoints are calibrated on the log target. Exponentiating the interval limits
produces asymmetric intervals in the original units used for scoring. The appendix NLL values for
these endpoints use the corresponding log-normal density in the original units.

Figures~\ref{fig:calibration-brier-endpoints} and
\ref{fig:calibration-interval-endpoints} show the full-grid mean score for each method and endpoint
before and after calibration. Each point averages 250 outer-fold evaluations. LightGBM and the
deterministic MLP supply only point predictions for regression, so their interval-score panels begin
after the constant residual scale is fitted.

\begin{figure*}[p]
  \centering
  \includegraphics[width=\textwidth]{calibration_brier_paper_blue.pdf}
  \caption{Full-grid mean Brier scores before and after temperature scaling for all 13 binary
  datasets and eleven pipelines. Open markers show uncalibrated scores and filled markers show
  calibrated scores. Lower is better.}
  \label{fig:calibration-brier-endpoints}
\end{figure*}

\begin{figure*}[p]
  \centering
  \includegraphics[width=\textwidth]{calibration_interval_paper_blue.pdf}
  \caption{Full-grid mean 90\% interval scores before and after method-specific calibration for
  all nine regression datasets. Open markers show native intervals and filled markers show
  calibrated intervals. LightGBM and deterministic MLP supply no native interval, so only their
  calibrated scores are shown. Axes are endpoint-specific; lower is better.}
  \label{fig:calibration-interval-endpoints}
\end{figure*}

\section{Global rank comparisons}
\label{app:global-rank-details}

The Friedman test finds differences among methods for all three evaluation measures
(Table~\ref{tab:global-rank-tests}).

\begin{table*}[!htbp]
\centering
\caption{Cross-dataset Friedman tests. $Q$ is the tie-corrected Friedman statistic.}
\label{tab:global-rank-tests}
\small
\begin{tabular}{@{}lrrr@{}}
\toprule
Evaluation measure & Datasets & $Q$ & $p$-value \\
\midrule
TDC leaderboard & 22 & 107.13 & $2.02\times10^{-18}$ \\
Calibrated Brier & 13 & 45.75 & $1.59\times10^{-6}$ \\
Calibrated interval score & 9 & 54.40 & $4.08\times10^{-8}$ \\
\bottomrule
\end{tabular}
\end{table*}

Table~\ref{tab:global-rank-holm} gives every mean rank and the Holm-adjusted comparison with ensemble
MC dropout. The adjustment includes all 55 pairwise comparisons within each evaluation measure.

\begin{table*}[!htbp]
\centering
\caption{Mean ranks and Holm-adjusted comparisons with ensemble MC dropout. Bold adjusted
$p$-values indicate a significant difference at the 5\% level. The adjustment accounts for all 55
pairwise method comparisons within each evaluation measure.}
\label{tab:global-rank-holm}
\scriptsize
\setlength{\tabcolsep}{4pt}
\begin{tabular}{@{}lrrrrrr@{}}
\toprule
& \multicolumn{2}{c}{TDC leaderboard} & \multicolumn{2}{c}{Brier} & \multicolumn{2}{c}{Interval score} \\
\cmidrule(lr){2-3}\cmidrule(lr){4-5}\cmidrule(l){6-7}
Method & Rank & $p_{\mathrm{Holm}}$ & Rank & $p_{\mathrm{Holm}}$ & Rank & $p_{\mathrm{Holm}}$ \\
\midrule
LightGBM & 6.36 & \textbf{0.0097} & 4.23 & 1.0000 & 3.56 & 1.0000 \\
Deterministic MLP & 8.70 & \textbf{$<0.0001$} & 8.38 & \textbf{0.0101} & 10.33 & \textbf{$<0.0001$} \\
MC dropout & 7.91 & \textbf{$<0.0001$} & 7.69 & 0.0690 & 5.67 & 1.0000 \\
SWAG & 7.32 & \textbf{0.0002} & 7.38 & 0.1431 & 7.78 & \textbf{0.0486} \\
Laplace & 8.16 & \textbf{$<0.0001$} & 9.23 & \textbf{0.0007} & 9.22 & \textbf{0.0014} \\
Post-hoc VBLL & 8.82 & \textbf{$<0.0001$} & 6.46 & 0.9364 & 8.33 & \textbf{0.0142} \\
Deep ensemble & 4.20 & 1.0000 & 4.54 & 1.0000 & 6.00 & 1.0000 \\
Ensemble MC dropout & 2.73 & -- & 3.54 & -- & 2.67 & -- \\
MultiSWAG & 3.86 & 1.0000 & 5.46 & 1.0000 & 5.00 & 1.0000 \\
MoLA & 3.20 & 1.0000 & 5.46 & 1.0000 & 3.44 & 1.0000 \\
Ensemble post-hoc VBLL & 4.73 & 1.0000 & 3.62 & 1.0000 & 4.00 & 1.0000 \\
\bottomrule
\end{tabular}
\end{table*}

\subsection{Computational scale}

Let $U=22\times10\times5\times5=5{,}500$ denote the number of
dataset--partition--model-seed--outer-fold units. Table~\ref{tab:compute-counts} counts shared
neural sweeps and ensemble backbones once, irrespective of how many UQ readouts they produce.

\begin{table}[!htbp]
\centering
\caption{Computation underlying the eleven-pipeline comparison.}
\label{tab:compute-counts}
\small
\begin{tabularx}{\linewidth}{@{}Xr@{}}
\toprule
Computation & Count \\
\midrule
LightGBM inner-CV fits & 880,000 \\
LightGBM outer refits & 5,500 \\
Shared MLP inner-CV fits & 880,000 \\
Single-model MLP outer refits & 5,500 \\
Ensemble-member MLP refits & 55,000 \\
\textbf{From-scratch model fits} & \textbf{1,826,000} \\
Checkpoint-resumed SWAG continuations & 60,500 \\
Post-hoc VBLL last-layer grid-point fits & 4,900,500 \\
Calibration applications & 60,500 \\
Outer-fold method evaluations & 60,500 \\
Score evaluations across the three main measures & 121,000 \\
\bottomrule
\end{tabularx}
\end{table}

The 55,000 ensemble-member entries are full MLP refits. The larger post-hoc VBLL count consists of
81 last-layer optimizations for each of 60,500 frozen feature maps, comprising 5,500 single models
and 55,000 ensemble members. Each optimization fits one prior-scale and Wishart-scale pair without
retraining the MLP backbone.
\FloatBarrier

\subsection{Additional evaluation measures}

Figures~\ref{fig:binary-metric-rails}--\ref{fig:regression-uq-metric-rails} show mean ranks for the
additional measures listed in Section~\ref{sec:methods}. Each score averages the same 250 outer-fold
values as the main analysis. The corresponding values on each dataset's original scale appear in
Figures~\ref{fig:appendix-native-binary-auroc}--\ref{fig:appendix-native-regression-coverage-error}.

\begin{figure*}[p]
  \centering
  \includegraphics[width=\textwidth,height=0.82\textheight,keepaspectratio]{appendix_binary_metric_rails.pdf}
  \caption{Mean ranks for five binary evaluation measures after temperature scaling. Each measure
  contains the same 13 datasets and eleven methods. The outlined marker identifies the lowest
  observed mean rank; coincident points have equal mean ranks.}
  \label{fig:binary-metric-rails}
\end{figure*}

\begin{figure*}[p]
  \centering
  \includegraphics[width=\textwidth]{appendix_regression_point_metric_rails.pdf}
  \caption{Mean ranks for regression point-prediction measures. Each measure contains the same nine
  datasets and eleven methods. The outlined marker identifies the lowest observed mean rank;
  coincident points have equal mean ranks.}
  \label{fig:regression-point-metric-rails}
\end{figure*}

\begin{figure*}[p]
  \centering
  \includegraphics[width=\textwidth]{appendix_regression_uq_metric_rails.pdf}
  \caption{Mean ranks for interval score, Gaussian NLL, and absolute coverage error after
  method-specific regression calibration. Each measure contains the same nine datasets and eleven
  methods. The outlined marker identifies the lowest observed mean rank; coincident points have
  equal mean ranks.}
  \label{fig:regression-uq-metric-rails}
\end{figure*}
\FloatBarrier
